\section{Introduction}
\label{sec:intro}

Containerization has become a cornerstone of cloud-native computing due to its efficiency and
scalability. Unlike traditional virtual machines, containers share the host OS kernel, leading to
reduced overhead but also introducing critical security risks through system calls
(syscalls)---the interface between user space and the kernel~\cite{ref1_jarkas_survey,ref3_sultan_survey}.

Syscalls expose a broad attack surface, making them attractive targets for privilege escalation,
container escape, and other exploits~\cite{ref5_combe_docker}. Existing solutions like seccomp
profiles offer only limited protection. Static profiles are hard to maintain and may block
legitimate operations or fail to detect sophisticated threats.

To overcome these limitations, we propose DeSFAM, Dynamic eBPF-Driven Syscall Filtering and
Anomaly Mitigation. DeSFAM provides real-time syscall profiling, deep anomaly detection, and
adaptive policy enforcement. It integrates:
\begin{enumerate}[(i)]
  \item Hybrid Syscall access listing using static and eBPF-based dynamic analysis;
  \item SyscallAD, a detection engine combining Variational Autoencoder (VAE) and Isolation
        Forest (iForest) for identifying anomalous syscall sequences; and
  \item Adaptive Policy Management that uses MITRE Adversarial Tactics, Techniques, and Common
        Knowledge (ATT\&CK) framework mappings, Common Vulnerabilities and Exposures (CVE)
        correlations, and eBPF LSM hooks for context-aware syscall control.
\end{enumerate}

DeSFAM enforces fine-grained syscall control using eBPF maps and LSM hooks, dynamically
blocking high-risk syscalls (e.g., \texttt{execve}, \texttt{ptrace}) and updating policies at
runtime without requiring container restarts. It combines a hybrid syscall profiling pipeline
with SyscallAD, a low-latency anomaly detector based on Variational Autoencoder (VAE) and
Isolation Forest. To prioritize threats, DeSFAM integrates contextual risk scoring by
correlating syscall behaviors with MITRE ATT\&CK tactics and known CVEs. Evaluations using the
DongTing dataset and real-world CVE scenarios demonstrate 94\% precision, 90\% recall,
sub-millisecond enforcement latency, and under 1\% performance overhead.

Our key contributions are:
\begin{itemize}
  \item A hybrid profiling pipeline for syscall allowlist generation.
  \item SyscallAD, a low-latency VAE + iForest anomaly detector.
  \item Contextual risk scoring using MITRE ATT\&CK and CVE correlations.
  \item Fine-grained syscall enforcement using eBPF maps and LSM hooks.
  \item A runtime feedback loop for policy auto-hardening based on violations.
\end{itemize}

The paper is organized as follows: Section~\ref{sec:background} covers syscall filtering and
eBPF fundamentals. Section~\ref{sec:related} discusses prior work. Section~\ref{sec:desfam}
describes DeSFAM's architecture. Section~\ref{sec:evaluation} presents evaluation results,
followed by discussion in Section~\ref{sec:discussion} and conclusion in
Section~\ref{sec:conclusion}.
